In the rational reconstruction that Cassirer broadly follows, Dalton's hypothesis gives rise to a fundamental difficulty. The empirical facts do not uniquely determine the relative atomic weights. Suppose hydrogen is taken as the unit: the same experimental composition can be represented using different atomic weights for oxygen, provided that the chemical formulas are altered accordingly. Thus, the experimental facts of composition alone cannot decide between, for example, $O=16$ and $O=8$. The same applies to sulfur, for which values such as $S=8, \quad S=16, \quad S=32, \ldots$ remain possible in the absence of further constraints \SF{274--75}{206--7}. In principle, chemists could have accepted this ambiguity as unavoidable; instead, they sought additional, independent relations that would allow them to determine the atomic weights unambiguously. Following \citets{Meyer1884} account, Cassirer shows how the \emph{univocal determination} of relative atomic weights emerges progressively, not from a single measurement, but from the mutual confirmation and correction of results obtained through independent physical laws. Cassirer, following Meyer, mentions three particularly important ones: Avogadro's rule\footnote{Equal volumes of gases at the same pressure and temperature contain equal numbers of molecules ($M \propto d$, where $M$ is molar mass and $d$ is gas density).}, relating molecular weights to gas volumes and vapor densities; the Dulong--Petit law, relating atomic weights to specific heats\footnote{The product of an element's atomic weight and its specific heat is approximately constant ($A \cdot c \approx 6.0\text{–}6.4$, with specific heat expressed in calories per gram per kelvin), within the law's range of applicability.}; and Mitscherlich's law of isomorphism, relating crystalline form to chemical composition and structure\footnote{Isomorphism provides evidence of chemical analogy and can therefore help establish corresponding chemical formulas.} \SF{275}{207}. No single rule can univocally determine all the relative atomic weights. Rather, these weights become increasingly determinate as several initially distinct series of phenomena are brought into mutual agreement: \qt{It is only the totality of all these different points of view, mutually confirming and correcting each other, that finally after many experiments gives a unitary table of atomic weights, and thus lays the basis of a definite system of chemical formulae}{Erst die Gesamtheit aller dieser verschiedenen Gesichtspunkte, die sich wechselseitig bestätigen oder korrigieren, ergibt schließlich nach mannigfachen Versuchen eine einheitliche Tabelle der Atomgewichte und legt damit den Grund zu einer eindeutigen chemischen Formelsprache} \SF{275}{207}.

Here the conditions of \q{scientific \ls{production}} must be distinguished from those of \qt{critical \ls{reflection}}{Die Bedingungen der wissenschaftlichen Produktion sind andere, als es die der kritischen Reflexion sind} \SF{279}{210}. A modern chemist might maintain that \qt{[t]he objective existence of the different types of atoms is presupposed; it is only necessary to discover their \ls{properties}, and to define them more exactly}{Die objektive Existenz der verschiedenen Arten von Atomen wird vorausgesetzt: nur ihre Eigenschaften gilt es noch zu ermitteln und quantitativ schärfer zu umgrenzen} \SF{275}{207}. Although scientists may regard the existence of atoms as their starting point, Cassirer, following Duhem and Ostwald, argues that a closer examination of the empirical grounds for their claims reveals precisely the opposite. The starting point consists of experimental laws, including the conservation of mass and the laws of definite, reciprocal, and multiple proportions, which establish invariant quantitative relations among the masses of substances involved in chemical reactions. In particular, the laws of definite and multiple proportions reveal the regularities governing the composition of chemical compounds, in contrast to the variable proportions characteristic of physical mixtures. The concept of \s{combining weight} is subsequently reinterpreted in terms of the \s{atomic weight} of discrete constituents, whose combinations account for the observed whole-number ratios. Yet the determination of these atomic weights requires the coordination of further independent empirical relations. Thus, \qt{It becomes clear that the atom is never the given \ls{starting-point}, but always only the \ls{goal} of our scientific statements}{Das Atom ist, wie jetzt sogleich deutlich wird, niemals der gegebene Ausgangspunkt, sondern immer nur der Endpunkt unserer wissenschaftlichen Aussagen} \SF{276}{208}.

